\subsection{Modules of fractions of graded modules}

Let A be a graded ring, M a graded A-module and \(\Delta\) the degree monoid; we shall assume in this no. that \(\Delta\) is a group. Recall (Algebra, Chapter II, § 11) that A and M are respectively direct sums of additive groups

\[
\mathbf {A} = \bigoplus_ {i \in \Delta} \mathbf {A} _ {i}, \quad \mathbf {M} = \bigoplus_ {i \in \Delta} \mathbf {M} _ {i}
\]

with \(\mathbf{A}_i\mathbf{A}_j\subset \mathbf{A}_{i + j}\) and \(\mathrm{A}_i\mathrm{M}_j\subset \mathrm{M}_{i + j}\) for all \(i,j\in \Delta\) . Let S be a multiplicative subset of A all of whose elements are homogeneous. For all \(i\in \Delta\) , we write \(\mathbf{S}_i = \mathbf{S}\cap \mathbf{A}_i\) and denote by \((\mathrm{S}^{-1}\mathbf{M})_i\) the set of elements \(m^{\prime}\) of \(\mathbf{S}^{-1}\mathbf{M}\) for which there exist elements \(j\) , \(k\) of \(\Delta\) , an element \(m\in \mathbf{M}_j\) and an element \(s\in \mathbf{S}_k\) such that \(j - k = i\) and \(m^{\prime} = m / s\) . If \((m_q')_{1\leqslant q\leqslant r}\) is a finite family of elements of \(\mathbf{S}^{-1}\mathbf{M}\) such that

\[
m _ {q} ^ {\prime} \in (\mathrm{S} ^ {- 1} \mathrm{M}) _ {i (q)},
\]

there exist elements \(j(q) \in \Delta\) and \(k \in \Delta\) , elements \(m_q \in \mathrm{M}_{j(q)}\) ( \(1 \leqslant q \leqslant r\) ) and \(s \in \mathrm{S}_k\) such that \(m_q' = m_q / s\) for \(1 \leqslant q \leqslant r\) (no. 1, Remark 2).

PROPOSITION 21. The ring \(\mathbf{S}^{-1}\mathbf{A}\) with the family \(((\mathbf{S}^{-1}\mathbf{A})_i)\) is a graded ring and \(\mathbf{S}^{-1}\mathbf{M}\) with the family \(((\mathbf{S}^{-1}\mathbf{M})_i)\) is a graded module over the graded ring \(\mathbf{S}^{-1}\mathbf{A}\) . The canonical mappings \(i_{\mathbf{A}}^{\mathbf{S}}\) and \(i_{\mathbf{M}}^{\mathbf{S}}\) are homogeneous of degree 0.

Let \(m \in M_j\) , \(s \in S_k\) , \(m' \in M_{j'}\) , \(s' \in S_{k'}\) and suppose that \(j - k = j' - k' = i\) ; then \((m/s) - (m'/s') = (s'm - sm')/ss'\) and \(s'm - sm' \in M_{j+k'} = M_{j'+k}\) and \(ss' \in S_{k+k'}\) , hence \((m/s) - (m'/s') \in (\mathrm{S}^{-1}\mathrm{M})_i\) by definition; this shows that the \((\mathrm{S}^{-1}\mathrm{M})_i\) are additive subgroups of \(S^{-1}\mathrm{M}\) . The sum of these groups is the whole of \(S^{-1}\mathrm{M}\) : for every \(x \in S^{-1}\mathrm{M}\) may be written as \(m/s\) where \(m \in \mathbf{M}\) , \(s \in \mathbf{S}\) ; \(s\) is homogeneous by hypothesis and \(m\) is a sum of homogeneous elements \(m_j\) ; hence \(x\) is the sum of the \(m_j / s\) , each of which belongs to a subgroup \((\mathbf{S}^{-1}\mathbf{M})_i\) . Finally the sum of the \((\mathbf{S}^{-1}\mathbf{M})_i\) is direct; for let us consider a finite family of elements \(x_q\) ( \(1 \leqslant q \leqslant n\) ) such that \(x_q \in (\mathbf{S}^{-1}\mathbf{M})_{i(q)}\) , where the indices \(i(q)\) are distinct, and suppose that \(\sum_{q=1}^{n} x_q = 0\) . Each \(x_q\) may be written as \(x_q = m_q / s\) where \(s \in S_k\) and \(m_q \in M_{i(q)+k}\) ; the hypothesis implies that there exists \(s' \in S\) such that \(s'\left(\sum_{q=1}^{n} m_q\right) = 0\) ; if the \(s'm_q\) were not all zero, we would have a contradiction since, if \(s' \in S_d\) , then \(s'm_q \in M_{i(q)+d}\) and the \(i(q) + d\) are all distinct. It follows that \(x_q = 0\) for every index \(q\) .

It is immediately verified that, if \(a \in (\mathrm{S}^{-1}\mathrm{A})_i\) and \(x \in (\mathrm{S}^{-1}\mathrm{M})_j\) , then \(ax \in (\mathrm{S}^{-1}\mathrm{M})_{i+j}\) . Applying this result to the case \(\mathbf{M} = \mathbf{A}\) , we see first that \(\mathrm{S}^{-1}\mathbf{A}\) is a ring graded by the \((\mathrm{S}^{-1}\mathrm{A})_i\) ; then we see that \(\mathrm{S}^{-1}\mathrm{M}\) is graded module over \(\mathrm{S}^{-1}\mathrm{A}\) . Finally, as \(1 \in \mathrm{A}_0\) , \(i_{\mathrm{A}}^{\mathrm{S}}\) and \(i_{\mathrm{M}}^{\mathrm{S}}\) are homogeneous of degree 0.

PROPOSITION 22. Let A (resp. B) be a graded ring of type \(\Delta\) , M (resp. N) a graded module over the graded ring A (resp. B), S (resp. T) a multiplicative subset of A (resp. B) all of whose elements are homogeneous, \(f: \mathrm{A} \to \mathrm{B}\) a homogeneous homomorphism of degree 0 such that \(f(\mathrm{S}) \subset \mathrm{T}\) and \(u: \mathrm{M} \to \mathrm{N}\) an A-linear mapping which is homogeneous of degree \(k\) . Then the homomorphism \(f': \mathrm{S}^{-1}\mathrm{A} \to \mathrm{T}^{-1}\mathrm{B}\) derived from \(f\) (no. 1, Proposition 2) is homogeneous of degree 0 and the \((\mathrm{S}^{-1}\mathrm{A})\) -linear mapping

\[
u ^ {\prime} \colon \mathrm{S} ^ {- 1} \mathrm{M} \rightarrow \mathrm{T} ^ {- 1} \mathrm{N}
\]

derived from f and u (no. 2, Proposition 5) is homogeneous of degree k.

This follows immediately from the equations \(f'(a / s) = f(a) / f(s)\) and \(u'(m / s) = u(m) / f(s)\) .

Finally we note that, if E is a graded A-algebra and S a multiplicative subset of A consisting of homogeneous elements, \(S^{-1}E\) with its \((S^{-1}A)\) -algebra structure (no. 8) and the graduation \((S^{-1}E)_{t}\) is a graded \(S^{-1}A\) -algebra, as follows immediately from the definitions.

